\section{Introduction}\label{sec:introduction}

The Whitney trick is the geometric step that turns algebraic cancellation
of intersections into disjoint embedded surfaces. In dimension four,
the Whitney discs needed for this step may themselves meet the surfaces
one is trying to separate. The disc embedding problem asks whether
algebraic intersection data, supplemented by suitable dual spheres,
suffice to produce the embedded discs. Its solution for good fundamental
groups underlies topological surgery and the five-dimensional
$s$-cobordism theorem for four-manifolds
\cite[Theorem~5.1A and Chapters~7 and~11]{FQ1990}.

Here is the formulation we use. Let $M$ be an oriented four-manifold,
let $\pi=\pi_1(M)$, and fix basing paths, or \emph{whiskers}, for all
surfaces. The equivariant intersection pairing $\lambda$ records an
intersection point with its sign and the element of $\pi$ obtained from
these paths; its values lie in $\Z[\pi]$. A family of sphere maps
$g_j$ is algebraically dual to disc maps $f_i$ if
$\lambda(f_i,g_j)=\delta_{ij}1$. A geometric dual meets its corresponding
disc in one transverse point and misses the other discs. We use the
reduced Wall invariant
\begin{equation}\label{eq:wall-convention}
 \widetilde\mu\colon\pi_2(M)\longrightarrow
 \Z[\pi]/\langle h-h^{-1},\,\Z\cdot1\rangle_{\mathrm{add}}.
\end{equation}
The quotient is additive, rather than a quotient by a group-ring ideal;
these are the oriented conventions of \cite[Section~1.3]{PRT2025}.

The disc embedding theorem applies to disc maps with disjoint locally
flatly embedded boundaries and algebraic dual spheres satisfying
$\lambda(g_i,g_j)=0=\widetilde\mu(g_i)$. When $\pi$ is good, it produces
disjoint locally flat discs with those boundaries, with geometric duals,
and with the induced boundary framings preserved for generically
immersed input discs \cite[Theorem~A]{PRT2025}.
The unrestricted disc embedding conjecture asks for this implication
without the good-group assumption. Theorem~\ref{thm:main} disproves it
by excluding every disjoint locally flat filling of the specified
boundary circles, even when no output framing or dual sphere is required.

\begin{theorem}\label{thm:main}
There exist a compact connected oriented smooth four-manifold $M$ with nonempty boundary, an integer $k\geq1$, proper generically immersed discs
\[
 f_i\colon(D^2,S^1)\longrightarrow(M,\partial M),\qquad 1\leq i\leq k,
\]
whose boundary restrictions form an embedding of $\coprod_iS^1$, and framed sphere maps $g_i\colon S^2\to M$ such that
\begin{equation}\label{eq:input-invariants}
 \lambda(f_i,g_j)=\delta_{ij}\,1,\qquad
 \lambda(g_i,g_j)=0,\qquad
 \widetilde\mu(g_i)=0
 \quad(1\leq i,j\leq k),
\end{equation}
but the given boundary circles do not bound pairwise disjoint proper locally flat embedded discs in $M$. Each $g_i$ can be chosen to be a smooth embedded sphere; different $g_i$ are allowed to intersect.
\end{theorem}

The boundary condition in this formulation is classical. The additional
force of the conclusion is its nonexistence assertion: changing the
relative homotopy classes of the prospective discs cannot evade the
obstruction. The vanishing assumptions are placed on the dual spheres;
no vanishing of the intersections among the input discs is needed.

\subsection{Good groups and the free-group problem}

Casson supplied the handle-embedding construction that makes the
simply connected theory possible. Freedman proved that every Casson
handle is homeomorphic, relative to its attaching region, to an open
two-handle \cite[Theorems~3.1 and~1.1]{Freedman1982}. Together these
results replace the unavailable smooth Whitney discs by topological
ones. Freedman--Quinn developed a version using capped gropes, obtained
by attaching surfaces along symplectic basis curves and closing the
upper curves with immersed caps \cite[Chapters~2--5]{FQ1990}.
The caps may intersect; the loops through these intersections carry
the fundamental-group information that must be controlled.

Goodness is a local condition on such models. In the formulation of
\cite[Definition~3.2]{PRT2025}, every height-$1.5$ disc-like capped
grope, equipped with a homomorphism from its fundamental group to the
group in question, must have an immersed disc in its four-dimensional
thickening with the same framed boundary whose double-point loops map
to the identity. A height-$1.5$ grope has a first surface stage with
further surface stages attached along one curve from each symplectic
pair; its remaining tips are capped.
The thickenings used in this test have free fundamental groups, which
explains the central role of the free-group case
\cite[pp.~510--511]{FT1995}.

Freedman--Teichner extended the positive theory to all groups of
subexponential growth \cite[Theorem~0.1]{FT1995}. Krushkal--Quinn gave
a different proof based on splitting gropes into dyadic branches and
organizing the resulting double-point words \cite[p.~408]{KQ2000}.
Their paper also contains the correction to the earlier height-raising
argument. The known class is closed under subgroups, quotients,
extensions and directed colimits; it includes finite and solvable groups as
well as groups of subexponential growth
\cite[Section~1.1]{PRT2025}. Thus exponential growth alone does not
mark the boundary of the theorem. The nonabelian free-group case
remained open in the 2025 account \cite{PRT2025}; see also
\cite[Chapter~23]{KOPR2021} for its place among four-manifold problems.

The construction of geometric dual spheres requires attention even
in the positive theory. Powell--Ray--Teichner supplied the missing
construction in the proof of Freedman--Quinn's theorem and also
controlled the homotopy classes of those duals
\cite[Theorem~A and Remark~1.3]{PRT2025}. Their correction concerns
output geometric duals and is distinct from the height-raising
correction in \cite{KQ2000}.

\begin{corollary}[Nongoodness of the free group]\label{cor:nongood}
The free group $F_2$ on two generators is not good in the sense of
Freedman--Quinn. Every group containing a subgroup isomorphic to $F_2$
is likewise not good.
\end{corollary}
\begin{proof}
The group $\pi_1(M)$ of the manifold in Theorem~\ref{thm:main} is
finitely presented because compact smooth manifolds have finite CW
homotopy type. It is not good: otherwise \cite[Theorem~A]{PRT2025},
applied to the discs and spheres in \eqref{eq:input-invariants}, would
supply the forbidden embedded discs.

Choose a finite generating list for $\pi_1(M)$ of length $r\geq2$,
padding if necessary. There is then a surjection
$F_r\twoheadrightarrow\pi_1(M)$. The elements $a^jba^{-j}$ for
$0\leq j<r$ freely generate a rank-$r$ subgroup of
$F_2=\langle a,b\rangle$, as follows from reduced-word normal form.
Subgroups and quotients of good groups are good
\cite[Section~1.1, p.~4]{PRT2025}. If $F_2$ were good, its subgroup
$F_r$ and then the quotient $\pi_1(M)$ would be good, a contradiction.
Finally, a good group containing a subgroup isomorphic to $F_2$ would
make $F_2$ good by subgroup closure, which proves the last assertion.
\end{proof}

The conclusion about $F_2$ is group-theoretic: it does not identify
$\pi_1(M)$ with $F_2$, and the geometric counterexample remains the
one in Theorem~\ref{thm:main}.

\subsection{Surgery and previous obstruction approaches}

The free-group problem also has concrete link formulations. Freedman
asked whether the Borromean rings are $A$--$B$ slice
\cite{Freedman1986AB}. Here each component is assigned a decomposition
$D^4=A_i\cup B_i$ extending the standard genus-one decomposition of
$S^3$. One seeks disjoint copies of all the pieces with attaching
curves given by the link and a parallel copy; each copy must extend
to an ambient self-homeomorphism of $D^4$. Freedman--Lin relate this
condition for the full generalized Borromean family to topological
surgery, equivalently to slicing the corresponding Whitehead doubles
with meridians freely generating the slice-complement group
\cite[pp.~91--92]{FL1989}.

Two results show why this obstruction program is delicate. Krushkal
proved that every link with vanishing pairwise linking numbers is
weakly $A$--$B$ slice, where the ambient-extension requirement is
removed \cite[Theorem~1 and Definition~2.4]{Krushkal2008}.
Freedman--Krushkal subsequently constructed homotopy $A$--$B$ slices
for the universal generalized Borromean family using the $2$-Engel
relation \cite[Theorem~1 and Definition~3.10]{FK2016}.
These are positive results for weakened embedding problems. The
present obstruction instead starts from the actual images of
hypothetical embedded discs and extracts restrictions on them through
representation deformations. The independent conjugates of all triple
commutators used here are specified in Section~\ref{sec:geometry};
they are not an invocation of the Engel relation.

Positive surgery results for free fundamental groups also persist in
important special cases. Krushkal--Lee solve the unobstructed surgery
problem for a degree-one normal map from a closed four-manifold to a
Poincare complex with free fundamental group, when the target
intersection pairing is extended from the integers
\cite[Theorem~1]{KL2002}. Their proof cuts along inverse images of
points under a map to a graph and reduces the stabilized problem to
the simply connected case. In our construction the corresponding
cuts have torus boundaries. Their boundary restriction data are
retained throughout the obstruction; the closed-manifold result does
not supply fillings for them \cite[Introduction, Remark]{KL2002}.

For precision, we use the following two unrestricted assertions in
the sense of \cite[Conjectures~1.5--1.6]{KPT2021}. The
\emph{topological surgery assertion} says that a degree-one normal
map from a compact four-manifold to a four-dimensional Poincare pair
$(Z,\partial Z)$, which is a $\Z[\pi_1(Z)]$-homology equivalence on the boundary and has
zero Wall surgery obstruction, is normally bordant relative to the
boundary to a homotopy equivalence. The \emph{$s$-cobordism assertion}
says that every compact five-dimensional topological $s$-cobordism
with a specified product on its boundary is homeomorphic to a product
extending that boundary structure. Neither assertion restricts the
fundamental group.

\begin{corollary}\label{cor:surgery-conjunction}
The unrestricted topological surgery and $s$-cobordism assertions
cannot both hold.
\end{corollary}
\begin{proof}
Together these assertions imply the unrestricted disc embedding
conjecture \cite[Section~3]{KPT2021}. Its algebraic hypotheses apply
to the finite family of discs and framed spheres in the compact
manifold of Theorem~\ref{thm:main}, so it
would give the excluded boundary filling.
\end{proof}

The corollary concerns their conjunction; it does not determine which
assertion fails. In particular, the free-group conclusion above is
not being used as a substitute for the quoted equivalence.

\subsection{The obstruction}
The construction starts with a smooth four-manifold having the homotopy type of a finite graph. A finite family of capped gropes inside it encodes every triple commutator of suitably conjugated end generators. Removing neighborhoods of the grope bases gives the actual manifold $M$. The shallow caps are the input discs. Product tori, compressed using the upper stages, give the spheres in \eqref{eq:input-invariants}. The proof checks their framings and all equivariant cancellation labels in $\pi_1(M)$.

Suppose disjoint locally flat discs with the prescribed boundaries existed. They would permit a graph map to be altered near their images and then cut along regular levels. The resulting three-dimensional walls have one torus boundary apiece, and the complementary four-dimensional region carries the triple commutator relations. An adjoint $\SL_2(\C)$ local system turns these relations into a cube-zero ideal in a representation deformation ring. A two-sheeted Mayer--Vietoris calculation supplies two complementary tangent spaces.

The comparison of representation deformations and flat connections has
a classical differential graded Lie algebra formulation in
Goldman--Millson \cite[Proposition~2.6 and Section~6]{GM1988}.
Here finite transition-cochain charts retain the equations themselves,
which will be needed with their full ideal structure.

To describe the resulting obstruction, label the walls $1,\ldots,d$.
Write $a_i\in\C^3$ for holonomy coordinates centered at the identity
along a chosen boundary circle on the $i$th wall's torus, and $y_i$ for its remaining deformation
coordinates. Its function $P_i(a_i,y_i)$ is a boundary-polarized version
of the Chern--Simons transgression integral
\cite[Section~3]{ChernSimons1974}.
Write $a=(a_i)_i$, and let $t$ be the remaining deformation parameters
of the cut four-manifold. Label its two copies of each wall so that
the chosen holonomy may vary on the positive copy and is trivial on
the negative copy. Their restrictions then have the form
$(a_i,Y_i^+(a,t))$ and $(0,Y_i^-(a,t))$.
Let $Y=(Y_i^+,Y_i^-)_i$, and let $C=(C_i^+,C_i^-)_i$ be independent
tuples with each $C_i^\pm$ of the same size as $y_i$. Put
\[
 \Psi(a,C)=\sum_i\bigl(P_i(a_i,C_i^+)-P_i(0,C_i^-)\bigr),
 \qquad
 v(a,t)=(\partial_C\Psi)(a,Y(a,t)).
\]
Here differentiation precedes restriction to the graph $C=Y(a,t)$.
If $\mathfrak r$ is the ideal of equations for the four-dimensional
representation chart, the two identities are
\[
                  (v)=\mathfrak r,
           \qquad \Psi(a,Y(a,t))\in\mathfrak r^2.
\]
The first identifies the whole ideal, including its nilpotent
information. If the derivatives missed an equation, the resulting first
obstruction to flatness would pair trivially with every wall variation.
Duality and the paired Mayer--Vietoris isomorphism would then kill its
curvature class, permitting a flat lift that contradicts the minimality
of the defect equations. The second identity comes from comparing the wall
connections on the closed boundary of the cut manifold: Stokes turns
their total value into an integral of curvature squared.

Choose one scalar coordinate $a_{i,1}$ on each wall and set
$p_i(y_i)=\partial_{a_{i,1}}P_i(0,y_i)$. The end-holonomy cube law
gives two distinct families of memberships, for every triple of labels:
\[
 a_{i,1}a_{j,1}a_{k,1}\in(v),\qquad
 p_i(Y_i^-)p_j(Y_j^-)p_k(Y_k^-)\in(v).
\]
The boundary term in polarized variation places each $p_i(Y_i^-)$
in the end-deviation ideal modulo $\mathfrak r$, giving the second family.
These are ideal memberships with coefficients, stronger than statements
about the common zero set.

The next step makes these identities available in positive
characteristic. All gradient jets of each original $P_i$ have
representatives in one fixed smooth complex algebraic chart. This
allows replacement of $P_i$, to sufficiently high order along its wall
ideal, by an algebraic germ. The replacement uses exactness and faithful
flatness of completion \cite[Tags 00MA--00MC]{Stacks} and differential
tests for symbolic powers \cite[Theorem~3.6 and Lemma~2.3]{DSGJ2020}.
The square-value identity then permits a
formal correction of $Y$ making $\Psi$ vanish exactly on the corrected
graph. An invertible change of gradient generators transports both
displayed cubic memberships to that graph. Only finitely many constants
and inverse Jacobians are needed for the new potentials, the correction
and its membership witnesses, so all these data specialize to a field
$\Bbbk$ of positive characteristic. The corrected graph need not describe
representations; the original end-holonomy cube law has already supplied
the two memberships that the remaining argument uses.

In characteristic $\ell$, quotienting by the separate $\ell$th powers
of the coordinates gives a finite-dimensional complex of differential
forms. The corrected graph at $a=0$ and its copy with the positive and
negative blocks exchanged have complementary tangent spaces, and the
potential vanishes on both. They determine cohomology classes with
nonzero pairing. Using this pairing and the cubic law for the $p_i$,
we construct, when there are at least five walls, two classes in a
deformation of the complex whose pairing is a unit in the deformation
ring. The cubic law for the $a_{i,1}$ then forces a nonzero parameter
product times that pairing to vanish. Section~\ref{sec:frobenius}
constructs this contradiction in the finite complex.

Lemma~\ref{alg:replacement} and Theorem~\ref{frob:no-data} are stated independently of the geometric construction. They are available for other problems meeting their respective stated hypotheses; in Theorem~\ref{frob:no-data} these include both cubic ideal-membership laws. No additional application is asserted here.

\subsection{Conventions and organization}
All manifolds used in the construction are compact and smooth until the hypothetical locally flat discs are introduced. Corners are rounded when needed. Boundary orientations use the outward-normal-first convention. The commutator is $[x,y]=xyx^{-1}y^{-1}$, and changes of basing are retained as explicit conjugating words. A framed surface has a trivialized real normal two-plane bundle. The symbol $\mathfrak m$ denotes the maximal ideal of the local parameter ring currently in use. Formal series are always completed at the specified reference representation or origin.

Section~\ref{sec:geometry} constructs the input discs and algebraic duals,
and converts a hypothetical boundary filling into smooth cut data.
Section~\ref{sec:deformation} obtains complementary restriction tangents
and the cubic holonomy law. Section~\ref{sec:connections} constructs the
wall potentials and proves the derivative-ideal and square-ideal
identities. Section~\ref{sec:algebra} replaces the potentials by algebraic
germs, corrects the restriction graph, and specializes the identities to
positive characteristic. Section~\ref{sec:frobenius} proves the finite
algebraic obstruction. Section~\ref{sec:conclusion} assembles these
implications to prove Theorem~\ref{thm:main}.
